% !TEX program = lualatex
% Transparencias de Fundamentos de las Matemáticas II
% Clase 08: Factorización en primos II

\documentclass[aspectratio=169,11pt]{beamer}

\usepackage{fontspec}
\usepackage[spanish,es-nodecimaldot]{babel}
\usepackage{amsmath,amssymb,mathtools}
\usepackage{booktabs}
\usepackage{csquotes}
\usepackage{xcolor}

\usetheme{Madrid}
\usecolortheme{default}
\setbeamertemplate{navigation symbols}{}
\setbeamertemplate{footline}[frame number]

\definecolor{FdMBlue}{RGB}{20,55,100}
\definecolor{FdMRed}{RGB}{130,30,30}
\definecolor{FdMGreen}{RGB}{25,100,60}

\setbeamercolor{structure}{fg=FdMBlue}
\setbeamercolor{title}{fg=white,bg=FdMBlue}
\setbeamercolor{frametitle}{fg=white,bg=FdMBlue}
\setbeamercolor{block title}{fg=white,bg=FdMBlue}
\setbeamercolor{block body}{bg=blue!3}
\setbeamercolor{alerted text}{fg=FdMRed}

\setmainfont{Latin Modern Roman}
\setsansfont{Latin Modern Sans}
\setmonofont{Latin Modern Mono}

\newcommand{\N}{\mathbb{N}}
\newcommand{\Z}{\mathbb{Z}}
\newcommand{\Q}{\mathbb{Q}}
\newcommand{\R}{\mathbb{R}}
\newcommand{\C}{\mathbb{C}}
\newcommand{\Pow}{\mathcal{P}}
\newcommand{\id}{\operatorname{id}}
\newcommand{\mcd}{\operatorname{mcd}}
\newcommand{\mcm}{\operatorname{mcm}}
\newcommand{\im}{\operatorname{Im}}

\title[FdM II -- Clase 08]{Factorización en primos II}
\subtitle{Unicidad y aplicaciones}
\author{Antonio Falcó Montesinos}
\institute{Universidad CEU Cardenal Herrera}
\date{Fundamentos de las Matemáticas II}

\begin{document}

\begin{frame}
  \titlepage
\end{frame}

\begin{frame}{Objetivos de la clase}
\begin{itemize}
  \item Enunciar la unicidad de la factorización.
  \item Aplicar factorización a \(\mcd\) y \(\mcm\).
  \item Probar irracionalidad de raíces sencillas.
\end{itemize}
\end{frame}

\begin{frame}{Mapa de la clase}
\tableofcontents
\end{frame}

\section{Unicidad}

\begin{frame}{Unicidad}
\begin{block}{Teorema Fundamental}
\begin{itemize}
  \item Todo entero \(n>1\) admite factorización en primos.
  \item Esa factorización es única salvo el orden de los factores.
  \item Se escribe de forma canónica agrupando potencias de primos.
\end{itemize}
\end{block}
\end{frame}

\begin{frame}{Unicidad}
\begin{block}{Lema clave}
\begin{itemize}
  \item Si \(p\) es primo y \(p\mid ab\), entonces \(p\mid a\) o \(p\mid b\).
  \item Este resultado conecta primalidad y divisibilidad.
  \item Es esencial para comparar dos factorizaciones.
\end{itemize}
\end{block}
\end{frame}

\begin{frame}{Unicidad}
\begin{block}{Comparar factorizaciones}
\begin{itemize}
  \item Si un primo divide un producto de primos, debe coincidir con alguno de ellos.
  \item Se cancelan factores uno a uno.
  \item Quedan los mismos primos con las mismas multiplicidades.
\end{itemize}
\end{block}
\end{frame}

\section{Aplicaciones}

\begin{frame}{Aplicaciones}
\begin{block}{\(\mcd\) y \(\mcm\)}
\begin{itemize}
  \item El \(\mcd\) toma exponentes mínimos de los primos comunes.
  \item El \(\mcm\) toma exponentes máximos de todos los primos presentes.
  \item Este método evita listar divisores.
\end{itemize}
\end{block}
\end{frame}

\begin{frame}{Aplicaciones}
\begin{block}{Irracionalidad}
\begin{itemize}
  \item Para probar que \(\sqrt{2}\) es irracional, se supone \(\sqrt{2}=a/b\) en forma irreducible.
  \item De \(a^2=2b^2\) se deduce que \(a\) es par.
  \item Entonces también \(b\) es par, contradiciendo la irreducibilidad.
\end{itemize}
\end{block}
\end{frame}

\section{Profundización}

\begin{frame}{Profundización}
\begin{block}{Ejemplo}
\begin{itemize}
  \item \(72=2^3\cdot 3^2\).
  \item \(180=2^2\cdot 3^2\cdot 5\).
  \item \(\mcd(72,180)=2^2\cdot 3^2=36\).
  \item \(\mcm(72,180)=2^3\cdot 3^2\cdot 5=360\).
\end{itemize}
\end{block}
\end{frame}

\begin{frame}{Profundización}
\begin{block}{Error frecuente}
\begin{itemize}
  \item Tomar exponentes máximos para el \(\mcd\).
  \item Usar solo primos comunes para el \(\mcm\).
  \item Olvidar multiplicidades en la factorización.
\end{itemize}
\end{block}
\end{frame}

\begin{frame}{Profundización}
\begin{block}{Pregunta}
\begin{itemize}
  \item ¿Qué propiedad se pierde si \(1\) fuese primo?
  \item ¿Por qué la unicidad es más fuerte que la existencia?
  \item ¿Cómo detectar si una fracción está en forma irreducible?
\end{itemize}
\end{block}
\end{frame}


\section{Detalle y práctica}

\begin{frame}{Detalle y práctica}
\begin{block}{Unicidad de la factorización}
\begin{itemize}
  \item La factorización prima existe y es única salvo el orden de los factores.
  \item La unicidad permite comparar enteros mediante exponentes de primos.
  \item De ahí se obtienen fórmulas para \(\mcd\) y \(\mcm\) usando mínimos y máximos de exponentes.
\end{itemize}
\end{block}
\end{frame}

\begin{frame}{Detalle y práctica}
\begin{block}{Actividad breve}
\begin{itemize}
  \item Factorizar \(360\) y \(840\).
  \item Calcular \(\mcd(360,840)\) y \(\mcm(360,840)\) desde los exponentes.
  \item Comprobar que \(\mcd(a,b)\mcm(a,b)=|ab|\) en este caso.
\end{itemize}
\end{block}
\end{frame}

\begin{frame}{Cierre}
\begin{itemize}
  \item La factorización en primos existe y es única.
  \item \(\mcd\) y \(\mcm\) se leen en los exponentes.
  \item La unicidad permite probar irracionalidad.
\end{itemize}
\end{frame}

\begin{frame}{Trabajo recomendado}
\begin{itemize}
  \item Releer las secciones correspondientes del manual.
  \item Identificar definiciones, símbolos nuevos y enunciados principales.
  \item Preparar dudas y ejemplos para el seminario asociado.
\end{itemize}
\end{frame}

\end{document}
